\documentclass[10pt,a4paper]{amsart}
\usepackage[margin=19mm]{geometry}
\usepackage{amsmath,amssymb,mathtools}
\usepackage[scr=rsfs,cal=euler]{mathalfa}
\usepackage{xfrac}
\usepackage[unicode,hidelinks,bookmarks=false]{hyperref}
\newcommand{\ie}{i.e.\ }
\newcommand{\cf}{cf.\ }
\newcommand{\resp}{respectively\ }
\newcommand{\Subsection}{\subsection}
\providecommand{\nobreakdash}{\nobreak\hbox{-}}
\setlength{\emergencystretch}{2em}
\multlinegap0pt
\usepackage{xcolor}
\usepackage{amssymb}
\usepackage{mathtools}
\usepackage{algorithm}\usepackage{algpseudocode}
\usepackage{tikz}
\usepackage{enumerate}
\newtheorem{assumption}{Assumption}
\newtheorem{definition}{Definition}
\newtheorem{theorem}{Theorem}
\def\Cij{\mathcal{C}_i (j)}
\def\Z{\field{Z}}%
\newcommand{\Zp}{\Z_{\geq 0}}
\def\co{\mathrm{co}}
\def\field#1{\mathbb #1}%
\def\ri{\mathrm{ri}}
\title{Distributed Multi-Step Model Predictive Control for Consensus}
\author{Navid Noroozi}
\date{Source: arXiv:2602.14714v2 (CC BY 4.0)}
\begin{document}
\maketitle

\noindent\emph{Source and licence.} Distributed Multi-Step Model Predictive Control for Consensus. Version arXiv:2602.14714v2 (CC BY 4.0), distributed by arXiv under the Creative Commons Attribution 4.0 International licence, http://creativecommons.org/licenses/by/4.0/, as recorded in the arXiv licence metadata for this exact version and shown on the abstract page https://arxiv.org/abs/2602.14714v2. Full licence notice: \texttt{licenses/arxiv-2602.14714-CC-BY-4.0.txt}.\par\smallskip

\noindent\textit{Editorial scope.} Counted unit: the article's theorem thm:moreau\_nes\_suf (source0.tex lines 480--499). The article's complete original proof of it is delivered with it (source0.tex lines 513--587, one \texttt{proof} environment of 75 lines, which the article places away from the statement). No mathematics is written, repaired or re-derived: the statement and the proof are the article's own lines, in the article's own notation, and no retained line is substituted or re-flowed.  Retained uncounted as the source-local context the proof needs: lines 284--286; lines 312--318; lines 401--403; lines 405--407; lines 437--444; lines 450--476.  Nothing was omitted from any retained span, so the package carries no omission record.  No retained line contains a citation, so the wrapper carries no bibliography and no bibliography entry.  No retained line contains a figure environment, an image-inclusion command or drawing code, so no image is delivered and the package declares no asset dependency.  Editorial formatting only, in addition: an AMS \texttt{amsart} preamble that restates the article's own declarations and macros used by the retained lines, namely the environments \texttt{assumption}, \texttt{definition}, \texttt{theorem} on separate counters, together with the standard packages those lines need.  The retained environments print in this document as Assumption 1, Definition 1, Theorem 1 in order of appearance, whereas the article prints the same items with the numbering of its own full document.  The complete native source of the article as distributed in the arXiv e-print for this exact version is bundled unchanged as \texttt{sources/194736/source0.tex}.
\begin{equation}\label{eq:agent_dyn}
x_i(t+1)=g_i(x_i(t),u_i(t)),\qquad i=1,\dots,\ell,
\end{equation}

\begin{assumption}[Joint spanning-tree connectivity]\label{ass:joint_tree}
There exists $B\in\mathbb{N}$ such that for every $j \in \Zp$, the union graph
\[
\mathcal{G}[j,j+B-1] := (\mathcal{V},\ \cup_{s=j}^{j+B-1}\mathcal{E}(s))
\]
contains a directed spanning tree. In the bidirectional case, this is equivalent to joint connectivity.
\end{assumption}

\begin{equation}\label{eq:H_def_moreau}
H(z(j)):=\operatorname{co}\{z_1(j),\dots,z_\ell(j)\}.
\end{equation}

\begin{equation}\label{eq:V_diameter}
V(z(j)):=\operatorname{diam}(H(z(j)))=\max_{p,q\in\mathcal{V}}\|z_p(j)-z_q(j)\|.
\end{equation}

\begin{definition}[Asymptotic consensus]\label{def:asymp_consensus}
Consider the diameter function $V$ as in~\eqref{eq:V_diameter} for a network of agents~\eqref{eq:agent_dyn} with an associated directed graph
$\mathcal{G}(j)=(\mathcal{V},\mathcal{E}(j))$.
The network achieves asymptotic consensus if 
\[
\limsup_{j\to\infty} V\big(z(j)\big) \to 0 .
\]
\end{definition}

\begin{assumption}[strict convexity]
\label{ass:asymp_moreau_A1}
Let $z(j)=(z_1(j),\dots,z_\ell(j))\in(\mathbb R^d)^\ell$ denote the agent variables at outer step $j\in\mathbb Z_{\ge0}$.
For each $i$ define the local hull
$
\Cij=\operatorname{co}\big(\{z_i(j)\}\cup\{z_k(j):k\in\mathcal N_i(j)\}\big).
$
There exists a constant $\eta\in[0,1)$ such that for every $i$ and $j$:
\begin{enumerate}
\item \textbf{convex-hull inclusion:}
\begin{equation}
z_i(j+1)\in \Cij .
\label{eq:A1a_eps_hull}
\end{equation}
\item \textbf{Relative interior (RI) strictness when disagreement is nonzero:}
If $z_i(j)$ is not equal to all neighbor values, i.e. if
\[
z_i(j)\notin \bigcap_{k\in\mathcal N_i(j)}\{z_k(j)\},
\]
then
\begin{equation}
z_i(j+1)\in \ri\big(\Cij\big).
\label{eq:A1b_ri_strictness}
\end{equation}
where $\mathrm{diam}(\mathcal C):=\sup\{\|x-y\|:x,y\in\mathcal C\}$ and $\partial \mathcal C$ is the boundary.
\end{enumerate}
\end{assumption}

\begin{theorem}[necessary and sufficient conditions for asymptotic consensus]
\label{thm:moreau_nes_suf}
Consider a network of agents~\eqref{eq:agent_dyn} with an associated directed graph $(\mathcal{V},\mathcal{E}(j))$.
Consider any trajectory $z(\cdot)$ generated by a distributed update law $u(\cdot)$.

\smallskip
\noindent\textbf{(Necessity)}  
Assume that asymptotic consensus holds for all initial conditions and for all trajectories of the update law,
and assume further that the update law is \emph{local} in the sense that $z_i(j+1)$ depends only on
$\{z_k(j):k\in\{i\}\cup\mathcal N_i(j)\}$.
Then necessarily:
\begin{enumerate}
\item[\textbf{(N1)}] the hull inclusion \eqref{eq:A1a_eps_hull} holds;
\item[\textbf{(N2)}] the joint spanning-tree condition of Assumption~\ref{ass:joint_tree} holds.
\end{enumerate}
\smallskip
\noindent\textbf{(Sufficiency)}  
If Assumptions~\ref{ass:asymp_moreau_A1} and \ref{ass:joint_tree} hold, then $z(\cdot)$ satisfies
asymptotic consensus in the sense of Definition~\ref{def:asymp_consensus}.
\end{theorem}

\begin{proof}[Proof of Theorem~\ref{thm:moreau_nes_suf}]
We prove necessity and sufficiency separately.

\paragraph{Necessity of (N1)}
Assume asymptotic consensus holds for all initial conditions and the update law is local.
Suppose by contradiction that \eqref{eq:A1a_eps_hull} fails for some $i$ and $j$, i.e.
$z_i(j+1)\notin \Cij$ for some admissible configuration $z(j)$.
Since $\Cij$ is convex and closed, by the separating-hyperplane theorem there exists a vector
$v\in\mathbb S^{d-1}$ such that
\[
\langle v, z_i(j+1)\rangle \;>\; \sup_{x\in \Cij} \langle v,x\rangle.
\]
Because $\Cij\subseteq H(z(j)):=\co\{z_1(j),\dots,z_\ell(j)\}$, the right-hand side is at most
$\sup_{x\in H(z(j))}\langle v,x\rangle = \max_k \langle v,z_k(j)\rangle$. Hence
\[
\max_k \langle v,z_k(j+1)\rangle \;\ge\; \langle v,z_i(j+1)\rangle
\;>\; \max_k \langle v,z_k(j)\rangle.
\]
Thus the support function of the global hull strictly increases in direction $v$, meaning the global hull
expands in that direction. By locality one can fix all agents outside $\{i\}\cup\mathcal N_i(j)$ so that this
expansion cannot be compensated by other agents at time $j+1$. This contradicts asymptotic consensus for all
initial conditions. Therefore \eqref{eq:A1a_eps_hull} is necessary.

\paragraph{Necessity of (N2)}
Assume by contradiction that Assumption~\ref{ass:joint_tree} fails. Then for every $B\in\mathbb N$ there exists
a time $j_B$ such that the union graph $\mathcal G[j_B,j_B+B-1]$ has no directed spanning tree. In particular,
there exists a nontrivial partition of the node set $V=S\cup T$ with $S,T\neq\emptyset$ such that no node in $T$
is reachable from $S$ over this union window. Choose initial conditions at time $j_B$ as two distinct constants:
$z_i(j_B)=a$ for $i\in S$, $z_i(j_B)=b$ for $i\in T$ with $a\neq b$. By locality, nodes in $T$ depend only on
states within their in-neighborhoods, and because $S$ cannot reach $T$ over the union window, the evolution of
$z_T$ over $[j_B,j_B+B]$ is independent of $z_S$. Consequently, one can select (or equivalently, construct) a
trajectory consistent with locality in which the two groups remain separated, contradicting consensus for all
initial conditions. Thus joint spanning-tree connectivity is necessary.

\paragraph{Sufficiency}
Assume Assumptions~\ref{ass:asymp_moreau_A1} and \ref{ass:joint_tree}.
Recal the global hull $H(z(j))$ as in~\eqref{eq:H_def_moreau} and its diameter $V(z(j))$ as in~\eqref{eq:V_diameter}.

\emph{Step 1 (nested hulls).} From \eqref{eq:A1a_eps_hull}, for every $i$ we have
$z_i(j+1)\in \Cij\subseteq H(z(j))$, hence $H(z(j+1))\subseteq H(z(j+1))$ and therefore $V(z(j+1))\le V(z(j))$.
Thus $V(z(j))$ is nonincreasing and converges to some limit $\limsup_{j \to \infty} \geq 0$.

\emph{Step 2 (leaving vertices under strictness).} Fix $j$ with $V(z(j))>0$.
Let $p$ be a vertex (extreme point) of $H(z(j))$ and let $I_p(j):=\{i: z_i(j)=p\}$ be the set of agents at that
vertex. Consider any $i\in I_p(j)$.
If all neighbor values equal $p$, then $\Cij=\{p\}$ and $z_i(j+1)=p$.
Otherwise $\Cij$ contains at least one point different from $p$, hence $p\in \partial \Cij$
and therefore $p\notin \ri(\Cij)$. By \eqref{eq:A1b_ri_strictness} we obtain
$z_i(j+1)\in \ri(\Cij)$, in particular $z_i(j+1)\neq p$. Thus an agent can stay at a vertex $p$ only
if all its current neighbor values equal $p$.

\emph{Step 3 (strict hull contraction over a connectivity window).}
Let $B$ be the window length from Assumption~\ref{ass:joint_tree}. Consider the union graph
$\mathcal G[j,j+B-1]$ which contains a directed spanning tree. Since $V(z(j))>0$, not all agents lie at the same
point. Hence there exists at least one vertex $p$ of $H(z(j))$ such that $I_p(j)\neq V$.
By the spanning-tree property, information from outside $I_p(\cdot)$ must enter $I_p(\cdot)$ somewhere over the
window; otherwise $I_p(\cdot)$ would be closed under in-neighbors across the union graph, preventing reachability
of all nodes. Hence there exists a time $\tau\in\{j,\dots,j+B-1\}$ and an edge $(k,i)\in \mathcal E(\tau)$ with
$i\in I_p(\tau)$ and $k\notin I_p(\tau)$. Then $i$ has a neighbor with value $\neq p$ and by Step 2 it must leave
$p$ at time $\tau+1$. Consequently $|I_p(\tau+1)|<|I_p(\tau)|$.

Repeating this argument over successive windows, one shows that if $V(z(j))>0$ then within at most $\ell B$ steps
(at most $\ell$ possible decrements of vertex-occupancy), at least one vertex of $H(z(j))$ loses all agents and
cannot remain an extreme point of the later hull. Therefore there exists $\hat B:=\ell B$ such that
\[
H(z(j+\hat B))\subsetneq H(z(j))\qquad\text{whenever }V(z(j))>0.
\]
Strict containment of compact convex sets implies strict decrease of the diameter, hence
$V(z(j+\hat B))<V(z(j))$ whenever $V(z(j))>0$.

\emph{Step 4 (conclude $\limsup_{j \to \infty} V(z(j))=0$).}
If $\limsup_{j \to \infty} V(z(j))>0$, then along the subsequence $j_m:=m\hat B$ we would have a strictly decreasing bounded-below
sequence $V(z(j_m))$, contradicting convergence to a positive limit. Hence $\lim_{j\to\infty}V(z(j))=0$.
This is asymptotic consensus by Definition~\ref{def:asymp_consensus}.
\end{proof}

\end{document}
